% =====================================================================
%  CHAPTER 11: 儿童少年伤害与暴力（对应大纲1第十一章）
%  参考PPT：相关伤害与暴力章节内容
% =====================================================================
\chapter{儿童少年伤害与暴力}
\setcounter{chapter}{11}
\chapterintro{
\textbf{掌握：}伤害的定义、分类及流行特征；意外伤害的Haddon模型与四E干预措施；道路交通伤害和溺水的预防；儿童虐待与忽视的分类及预防；自杀与自伤的流行特征及干预；校园暴力与网络暴力的预防。\newline
\textbf{熟悉：}非自杀性自伤（NSSI）的识别与干预；校园欺凌的特征与预防。
}

% =====================================================================
%  11.1 儿童少年意外伤害
% =====================================================================
\section{儿童少年意外伤害}

\subsection{概述}

\begin{defbox}
\textbf{伤害（injury）：}因机械能、热能、电能、化学能等能量传递或干扰超过人体组织的耐受水平而造成的组织损伤，或由于窒息导致的缺氧而引起的损害。

\textbf{意外伤害（unintentional injury）：}非主观故意行为导致的伤害，即无目的性、无意识的伤害。主要包括交通事故、溺水、跌落、烧伤/烫伤、中毒等。

\textbf{伤害与"意外"的区别：}伤害是一类疾病，在国际疾病分类（ICD-10）中已有明确分类。意外伤害并非"意外"——它是有因可寻、有规律可循、可预测、可预防的公共卫生问题。
\end{defbox}

\subsection{流行特征和危险因素}

\begin{keybox}
\textbf{一、全球流行特征}
\begin{enumerate}[leftmargin=*]
    \item \imp{伤害是儿童少年的第一位死因}：全球每年约95万18岁以下儿童少年死于伤害，其中意外伤害占90\%以上。
    \item \imp{五大伤害死因（全球儿童青少年）}：
    \begin{enumerate}
        \item 道路交通伤害
        \item 自杀（故意伤害）
        \item 跌落
        \item 人际间暴力
        \item 溺水
    \end{enumerate}
    \item 全球儿童伤害死亡中，95\%发生在中低收入国家。
\end{enumerate}

\textbf{二、分布特征}
\begin{enumerate}[leftmargin=*]
    \item \imp{地区分布}：发展中国家伤害死亡率显著高于发达国家；农村高于城市。
    \item \imp{人群分布}
    \begin{itemize}
        \item 年龄：不同年龄阶段高发伤害类型不同——婴幼儿以跌落、烧烫伤、窒息为主；学龄儿童以溺水、交通事故为主；青少年以交通事故、暴力为主。
        \item 性别：男性伤害发生率普遍高于女性，男女比约1.5:1--3:1。
        \item 种族/民族：不同种族间存在差异，与社会经济状况相关联。
    \end{itemize}
    \item \imp{地点分布}：不同伤害类型发生地点不同——家庭（跌落、烧烫伤）、学校（运动伤害）、道路（交通事故）、水域（溺水）。
    \item \imp{时间分布}：交通事故和溺水有明显季节性——夏季（7--8月）为高发期。
\end{enumerate}

\textbf{三、危险因素}
\begin{enumerate}[leftmargin=*]
    \item \imp{宿主因素（host factors）}
    \begin{itemize}
        \item 年龄：不同发育阶段面临不同的伤害风险（婴幼儿运动能力发展不协调、学龄儿童探索行为增多、青少年冒险行为增加）。
        \item 性别：男性冒险倾向更高，暴露于危险环境的机会更多。
        \item 种族：部分种族因社会经济地位等因素，伤害风险较高。
        \item 心理行为因素：冲动性、攻击性、注意力不集中、冒险行为等增加伤害风险。
    \end{itemize}
    \item \imp{环境因素}
    \begin{itemize}
        \item 自然环境：气候（高温/暴雨/冰雪）、地形（水域/山体）、自然灾害。
        \item 社会环境：经济发展水平、社区安全设施、家庭监护状况、学校安全教育等。
    \end{itemize}
\end{enumerate}
\end{keybox}

\begin{tipbox}
\textbf{伤害的流行病学意义：}伤害已成为全球儿童青少年死亡和残疾的最主要原因。我国0--14岁儿童中，意外伤害是第一死亡原因，每年有近5万名儿童因伤害死亡。伤害造成的疾病负担（DALYs）巨大，但通过有效的公共卫生干预，大部分伤害是可以预防的。
\end{tipbox}

\subsection{伤害预防控制理论}

\subsubsection{Haddon模型}

\begin{keybox}
\textbf{Haddon模型（Haddon Matrix）：}由William Haddon于1970年提出，是伤害预防控制最经典的理论框架。该模型将伤害事件分解为三个因素（宿主、致伤因子、环境）和三个阶段（事件前、事件中、事件后），构成$3\times3$矩阵。

\textbf{三因素（3 factors）：}
\begin{enumerate}[leftmargin=*]
    \item \imp{宿主（host）}：受害者本身——年龄、性别、行为、知识、技能等。
    \item \imp{致伤因子（agent）}：导致伤害的能量或物体——车辆、水、火、化学品等。
    \item \imp{环境（environment）}：
    \begin{itemize}
        \item 物理环境：道路状况、水域安全设施、建筑安全等。
        \item 社会环境：法律法规、社会规范、经济水平等。
    \end{itemize}
\end{enumerate}

\textbf{三阶段（3 phases）：}
\begin{enumerate}[leftmargin=*]
    \item \imp{事件前阶段（pre-event phase）}：伤害发生之前——预防伤害的发生（一级预防）。
    \item \imp{事件中阶段（event phase）}：伤害正在发生之际——减轻伤害的严重程度（二级预防）。
    \item \imp{事件后阶段（post-event phase）}：伤害发生之后——降低伤残和死亡率（三级预防）。
\end{enumerate}

\textbf{示例——道路交通伤害的Haddon模型：}
\begin{table}[H]
\centering
\caption{道路交通伤害的Haddon矩阵}
\begin{tabularx}{\textwidth}{l>{\raggedright\arraybackslash}p{3.5cm}>{\raggedright\arraybackslash}p{3.5cm}>{\raggedright\arraybackslash}p{3.5cm}}
\hline
& \textbf{宿主} & \textbf{致伤因子} & \textbf{物理-社会环境} \\
\hline
\textbf{事件前} & 安全教育、驾驶技能 & 车辆安全性能、制动系统 & 道路设计、交通法规 \\
\hline
\textbf{事件中} & 是否使用安全带/头盔 & 车辆碰撞保护能力 & 道路护栏、安全气囊 \\
\hline
\textbf{事件后} & 急救知识、健康基础 & 车辆自燃风险 & 急救系统、医院条件 \\
\hline
\end{tabularx}
\end{table}
\end{keybox}

\subsubsection{四E干预}

\begin{keybox}
\textbf{四E干预措施：}基于Haddon模型的伤害预防综合干预策略。

\begin{enumerate}[leftmargin=*]
    \item \imp{教育干预（Education）}：通过健康教育提高公众对伤害的认识，增强安全意识和自我保护能力。包括：学校安全教育课程、社区宣传活动、驾驶员安全教育、家长安全知识培训等。

    \item \imp{工程干预（Engineering）}：通过改进产品和环境设计来减少伤害风险。包括：安全座椅和头盔设计、车辆安全性能改进、道路安全设施建设、儿童安全包装、游泳池围栏、防火建筑材料等。

    \item \imp{经济干预（Economic）}：通过经济手段影响行为选择，促进安全决策。包括：对安全产品的补贴或税收优惠、对危险产品的税收或处罚、保险激励机制等。经济干预也被称为强制干预。

    \item \imp{强制干预（Enforcement）}：通过法律法规强制实施安全措施。包括：交通法规（限速、酒驾禁令）、安全生产法规、儿童安全座椅法规、学校安全管理制度等。
\end{enumerate}

\textbf{四E干预的特点：}四E干预措施相互补充，综合运用效果更佳。单独依靠教育干预效果有限，必须同时配合工程改进、经济激励和法律强制，才能实现伤害的有效预防。
\end{keybox}

\subsection{道路交通伤害}

\begin{keybox}
\textbf{一、概述}
\begin{enumerate}[leftmargin=*]
    \item \imp{全球状况}：道路交通伤害是全球10--24岁人群的第一位死因，每年约119万人死于道路交通事故。
    \item \imp{中国状况}：道路交通伤害是中国0--14岁儿童的第二位伤害死因（仅次于溺水），占儿童伤害死亡的较大比例。
    \item \imp{发生特点}：
    \begin{itemize}
        \item 高发季节：夏季（7--8月）最为集中
        \item 高发时段：下午放学高峰时段
        \item 易受伤部位：头部损伤最为常见
    \end{itemize}
\end{enumerate}

\textbf{二、危险因素}
\begin{enumerate}[leftmargin=*]
    \item \imp{个体因素}：年龄（儿童判断力不足、注意力易分散）、性别（男生发生率高于女生）、冒险行为（闯红灯、不佩戴头盔）、安全意识薄弱。
    \item \imp{家庭因素}：家长监护不足、家长安全意识淡薄、家庭社会经济地位低下。
    \item \imp{交通环境因素}：道路安全设施不完善（缺乏人行道、信号灯、减速带）、交通流量大、车辆速度快、混行交通。
\end{enumerate}

\textbf{三、预防措施}
\begin{enumerate}[leftmargin=*]
    \item \imp{健康教育}：学校交通安全教育课程；培养安全过马路、安全骑行的行为习惯；提高危险识别和自我保护能力。
    \item \imp{环境干预}：优化学校周边交通环境；设置人行横道、信号灯、减速带；建立安全上学路线。
    \item \imp{立法与执法}：强制使用儿童安全座椅和安全带；严格限速法规；酒驾零容忍政策；加强交通执法力度。
    \item \imp{工程干预}：改进车辆安全设计（儿童安全座椅、安全气囊）；头盔设计优化；道路基础设施安全改造。
\end{enumerate}
\end{keybox}

\begin{exambox}
\textbf{易考：}
\begin{enumerate}[leftmargin=*]
    \item Haddon模型的三因素、三阶段及其在道路交通伤害中的应用
    \item 四E干预措施的具体内容及之间的相互关系
    \item 道路交通伤害的危险因素及综合预防策略
\end{enumerate}
\end{exambox}

\subsection{溺水}

\begin{keybox}
\textbf{一、概述}
\begin{enumerate}[leftmargin=*]
    \item \imp{全球状况}：溺水是全球5--14岁儿童意外伤害死亡的第二位原因，每年约有17.5万儿童死于溺水。
    \item \imp{中国状况}：溺水是中国儿童意外伤害死亡的第一位原因，约占儿童伤害死亡的40\%，是威胁我国儿童生命的"头号杀手"。
    \item \imp{发生特点}：男性溺水率高于女性；农村高于城市；主要发生于自然水域（河流、湖泊、水库、池塘）；高发季节为夏季（6--8月）。
\end{enumerate}

\textbf{二、危险因素}
\begin{enumerate}[leftmargin=*]
    \item \imp{个体因素}：年龄（5--14岁高发）、男性、不会游泳或水性不佳、对危险水域认知不足、冒险和好奇心强烈。
    \item \imp{危险水域}：开放性水域缺乏安全设施（护栏、警示牌、救生设备）；水域环境复杂（暗流、水深、水温低）。
    \item \imp{监管缺乏}：家长或监护人疏于看管，儿童独自或结伴到水域玩耍。
    \item \imp{同伴压力}：同伴怂恿下水、相互模仿冒险行为，群体情境下风险意识降低。
\end{enumerate}

\textbf{三、预防措施}
\begin{enumerate}[leftmargin=*]
    \item \imp{管理干预}：
    \begin{itemize}
        \item 对开放性危险水域设置护栏、警示标志
        \item 配备救生设备和专职救生人员
        \item 填埋或改造危险积水坑、废弃水塘
        \item 建立重点水域巡查和报告制度
    \end{itemize}
    \item \imp{健康教育}：
    \begin{itemize}
        \item 学校防溺水安全教育专题课程
        \item "六不一会"教育：不私自下水游泳、不擅自与他人结伴游泳、不在无家长或教师带领的情况下游泳、不到无安全设施和无救援人员的水域游泳、不到不熟悉的水域游泳、不熟悉水性的学生不擅自下水施救；学会基本的自护自救方法
        \item 家长防溺水安全知识宣传
    \end{itemize}
    \item \imp{游泳技能培养}：推广学校游泳课程，提高儿童少年的游泳能力和水中自救技能。但需注意：掌握游泳技能不能完全替代安全防护措施，需与其他措施配合使用。
\end{enumerate}
\end{keybox}

% =====================================================================
%  11.2 儿童少年故意伤害
% =====================================================================
\section{儿童少年故意伤害}

\begin{defbox}
\textbf{故意伤害（Intentional Injury）：}是指有计划、有目的地加害自己和他人的伤害行为，主要包括自杀和他杀、自伤、虐待/疏忽、家庭/社会暴力、强奸以及与毒品/酒精有关的伤害等。与意外伤害（无目的性）相比，故意伤害具有明确的目的性和主观意图。
\end{defbox}

\subsection{儿童虐待与忽视}

\begin{keybox}
\textbf{一、定义与分类}
\begin{enumerate}[leftmargin=*]
    \item \imp{WHO定义}：儿童虐待（child maltreatment）是指对18岁以下儿童实施的躯体虐待（或/和）情感虐待、性虐待、忽视及经济剥削等，对儿童的健康、生存、发展或尊严造成实际或潜在伤害的一切行为。

    \item \imp{四种类型：}
    \begin{enumerate}
        \item 躯体虐待（physical abuse）：蓄意对儿童造成躯体伤害或痛苦。
        \item 情感虐待（emotional abuse）：持续地对儿童进行不适当的对待，包括贬低、威胁、恐吓、歧视或其他非躯体的敌意行为，对儿童的心理和情感发育造成损害。
        \item 性虐待（sexual abuse）：使尚未发育成熟的儿童参与其不能完全理解、无法表达知情同意的，或触犯法律、违反社会禁忌的性活动。
        \item 忽视（neglect）：父母或其他照护者未能提供儿童生长发育所必需的基本条件（食物、衣物、住所、医疗、教育、安全监护和情感关怀等），从而对儿童健康或发育造成损害或潜在损害。
    \end{enumerate}
\end{enumerate}

\textbf{二、中国流行状况}
\begin{enumerate}[leftmargin=*]
    \item 儿童躯体虐待发生率约20.3\%。
    \item 儿童性虐待发生率约18.2\%。
    \item 儿童忽视（包括躯体忽视、情感忽视、教育忽视、医疗忽视等）发生率约26.0\%。
    \item 虐待和忽视常同时存在，多种类型重叠发生。
\end{enumerate}

\textbf{三、危险因素}
\begin{enumerate}[leftmargin=*]
    \item \imp{个体层面}：儿童年龄（年幼儿童更易受虐待）、性别、特殊需求儿童（残疾、智力障碍）。
    \item \imp{家庭层面}：父母精神健康问题（抑郁、物质滥用）；父母有虐待史；家庭暴力；家庭功能失调；单亲家庭或社会经济地位低下。
    \item \imp{社会层面}：社会经济不平等；社会对暴力和体罚的容忍；儿童保护法律体系不完善；社区支持和监督机制薄弱。
    \item \imp{文化层面}："不打不成器"等传统观念；对儿童权利认识不足。
\end{enumerate}

\textbf{四、预防措施}
\begin{enumerate}[leftmargin=*]
    \item \imp{法律法规建设}：完善儿童保护法律法规体系；强制报告制度（医务人员、教师、社会工作者等发现儿童虐待需依法报告）；建立儿童虐待案件的处理和追责机制。
    \item \imp{健康教育}：面向全社会的反儿童虐待宣传教育；普及儿童权利知识；改变"体罚合理"等传统观念。
    \item \imp{家庭支持}：提供家庭教育指导服务；对有风险的家庭提供经济援助和心理支持；开展积极育儿技能培训。
    \item \imp{学校-社区合作}：建立学校与社区联动的儿童保护网络；学校开展预防性虐待和自我保护教育；发挥社区监督和支持功能。
    \item \imp{监测系统}：建立和完善儿童虐待与忽视的监测报告系统，及时识别、报告和干预虐待事件。
\end{enumerate}
\end{keybox}

\subsection{自杀与自伤}

\begin{keybox}
\textbf{一、概述}
\begin{enumerate}[leftmargin=*]
    \item \imp{定义}：自杀（suicide）是个体在意识清醒的情况下，自愿而非被强迫地采用残害自己生命的手段结束自己生命的行为。它是具有严重社会不良后果的，全球性的重大公共卫生问题。
    \item \imp{流行病学}：自杀是全球15--29岁人群的第四大死因。每年全球约有70万人死于自杀。
    \item \imp{三个层次：}
    \begin{enumerate}
        \item 自杀意念（suicidal ideation）：有结束自己生命的想法，但尚未付诸行动。
        \item 自杀未遂（attempted suicide）：采取了伤害自己的行为但未导致死亡。
        \item 自杀死亡（completed suicide）：以死亡为结局的自杀行为。
    \end{enumerate}
\end{enumerate}

\textbf{二、青少年自杀流行特征}
\begin{enumerate}[leftmargin=*]
    \item \imp{全球状况}：12--15岁青少年中，自杀意念报告率约16.5\%，自杀计划率约16.5\%，自杀未遂率约16.4\%。
    \item \imp{中国学生群体}：自杀意念报告率男性约16.9\%，女性约24.6\%（女性高于男性）；自杀未遂率约2--5\%。中国青少年自杀未遂报告率相对较低，但自杀是全球15--29岁人群的第四大死因。
    \item \imp{性别差异}：女性自杀意念和自杀未遂率高于男性，但男性自杀死亡率高于女性（男性手段更致命）。
\end{enumerate}

\textbf{三、主要解释模型}
\begin{enumerate}[leftmargin=*]
    \item \imp{应激-素质模型（stress-diathesis model）}：自杀是应激因素（生活压力事件）与个体素质（易感性）交互作用的结果。当个体的素质脆弱性（如冲动性、悲观、神经生物学因素）与应激事件相遇时，自杀风险显著增加。

    \item \imp{多因素模型（multi-factor model）}：自杀是生物、心理、社会等多层面因素综合作用的结果，包括遗传因素、精神障碍（抑郁症等）、人格特质、家庭环境、社会支持、文化因素等。
\end{enumerate}

\textbf{四、自杀的警告信号与具体表现}

\textbf{警告信号（warning signs）：}
\begin{enumerate}[leftmargin=*]
    \item 言语信号：表达"活着没意思"、"我不想活了"、"希望自己从没出生过"等。
    \item 行为信号：将自己的物品送人、写遗书、突然的社会退缩、自伤行为、寻求自杀手段等。
    \item 情绪信号：持续的绝望感、情绪极度低落或突发的平静（已做决定后的平静）、焦虑不安、愤怒或冲动。
    \item 情境信号：经历重大丧失（亲人死亡、关系破裂）、学业严重受挫、被欺凌/虐待等。
\end{enumerate}

\textbf{自杀的具体表现（五类）：}
\begin{enumerate}[leftmargin=*]
    \item \imp{语言方面}：话里透露想死的念头，或在作文、诗词、周记中有所表现。
    \item \imp{行为方面}：日常行为习惯变化大，与以往行为模式显著不同。
    \item \imp{环境应付困难}：常有家庭重大变故、重要人际关系结束、失恋、升学考试失败等易导致产生绝望感和自杀的事件。
    \item \imp{心理行为异常}：退出现有人际交往，与世隔绝，强烈孤独感等。
    \item \imp{抑郁症表现}：绝望、低自尊、无助感、对许多事物失去兴趣等，对自杀行为有较强的预示作用。
\end{enumerate}

\textbf{五、预防干预——三级预防策略}
\begin{enumerate}[leftmargin=*]
    \item \imp{一级预防（通用性预防）}：面向全体人群，旨在减少自杀的危险因素，增强保护因素。包括：学校心理健康教育、生命教育、情绪管理和应对技能培训、营造支持性校园环境、限制自杀手段的可及性（如农药管理、高楼防护）。

    \item \imp{二级预防（选择性预防）}：针对高危人群（有自杀意念或自伤行为的学生），进行早期识别和干预。包括：学校心理筛查和风险评估、心理咨询和心理疏导、危机热线的建立和推广、对教师和同伴进行"守门人"培训。

    \item \imp{三级预防（针对性预防）}：针对自杀未遂者，防止再次自杀。包括：心理治疗和医学康复、长期跟踪随访、社会支持系统的重建、回归学校和社会后的关怀。
\end{enumerate}
\end{keybox}

\subsection{非自杀性自伤（NSSI）}

\begin{keybox}
\textbf{一、定义}
\begin{enumerate}[leftmargin=*]
    \item \imp{非自杀性自伤（Non-Suicidal Self-Injury, NSSI）}：不以死亡为目的，故意、直接地对自身身体组织造成伤害的行为，且不被社会文化所认可。常见方式包括割伤、划伤、烧伤、撞击、咬伤等。
    \item \imp{与自杀的区别}：NSSI不以结束生命为目的，而是以缓解负面情绪、表达痛苦、自我惩罚等为目的。但NSSI是未来自杀行为的重要预测因素。
\end{enumerate}

\textbf{二、流行特征}
\begin{enumerate}[leftmargin=*]
    \item 全球青少年NSSI发生率约17.2\%。
    \item 中国青少年NSSI发生率约12.2--17.2\%。
    \item 女性发生率高于男性。
    \item 高发年龄：12--18岁（青春期至青年早期）。
\end{enumerate}

\textbf{三、高危因素}
\begin{enumerate}[leftmargin=*]
    \item \imp{个体层面}：情绪调节困难、冲动性人格、低自尊、童年期虐待经历、精神障碍（抑郁、焦虑、边缘性人格障碍）。
    \item \imp{家庭层面}：家庭冲突、亲子沟通不良、父母忽视或过度控制、家庭暴力史。
    \item \imp{学校层面}：学业压力过大、校园欺凌受害经历、同伴中有NSSI行为（模仿/传染效应）、师生关系紧张。
\end{enumerate}

\textbf{四、干预措施}
\begin{enumerate}[leftmargin=*]
    \item \imp{通用性干预（一）}：面向全体学生的心理健康教育；情绪调节和压力管理技能培训；营造积极的学校氛围。
    \item \imp{心理治疗（二）}：认知行为治疗（CBT）、辩证行为治疗（DBT）对NSSI有较好效果；帮助青少年学习替代性应对策略（如分散注意力、表达感受、放松训练等）。
    \item \imp{药物治疗（三）}：针对共患精神障碍（如抑郁症、焦虑症）使用相应药物，如SSRIs类药物。
\end{enumerate}
\end{keybox}

% =====================================================================
%  11.3 校园暴力与网络暴力
% =====================================================================
\section{校园暴力与网络暴力}

\subsection{校园暴力}

\begin{keybox}
\textbf{一、定义与类型}
\begin{enumerate}[leftmargin=*]
    \item \imp{定义}：校园暴力（school violence）是指在学校相关环境中，针对学生、教师或学校设施进行的暴力行为。包括发生在学校内以及上下学途中的暴力事件。
    \item \imp{主要类型}
    \begin{enumerate}
        \item 躯体暴力（physical violence）：打、踢、推搡、使用武器等造成身体伤害的行为。
        \item 言语/情感暴力（verbal/emotional violence）：辱骂、嘲笑、威胁、恐吓、孤立、散布谣言等造成心理伤害的行为。
    \end{enumerate}
\end{enumerate}

\textbf{二、表现形式}
\begin{enumerate}[leftmargin=*]
    \item \imp{学生间暴力}：最普遍的校园暴力形式，包括校园欺凌（欺凌者、受害者、旁观者）、打架斗殴、群体暴力等。
    \item \imp{师生间暴力}：教师对学生的体罚或侮辱、学生对教师的攻击行为。
    \item \imp{校外人员介入}：校外人员进入校园实施的暴力行为。
\end{enumerate}

\textbf{三、三级预防策略}
\begin{enumerate}[leftmargin=*]
    \item \imp{一级预防}：面向全体学生，创建安全和谐的校园环境。包括：学校反暴力政策和制度建设；校园安全设施和环境改造；生命教育和反暴力教育；社交-情绪学习（SEL）项目；积极的学校氛围营造。
    \item \imp{二级预防}：针对高风险学生群体（有暴力倾向或受害风险者）进行早期识别和干预。包括：心理行为筛查；愤怒管理和冲突解决技能训练；心理咨询和小组辅导；家长干预和家庭教育指导。
    \item \imp{三级预防}：对已发生暴力事件的处理和善后。包括：暴力事件的及时制止和妥善处置；对受害者的心理援助和医学救治；对施暴者的教育和行为矫正；防止再次发生。
\end{enumerate}
\end{keybox}

\begin{defbox}
\textbf{校园欺凌（school bullying）：}是校园暴力的一种常见形式，指一个或多个学生反复、持续地以负面行为对待另一个学生，且双方力量存在不平衡。特征是：蓄意性、反复性、力量不平衡。形式包括：直接欺凌（打、踢、推搡）和间接欺凌（孤立、散布谣言）。
\end{defbox}

\begin{keybox}
\textbf{四、校园暴力的流行特征与危害}

\textbf{1. 冰山模式（Iceberg Model）：}
暴力伤害符合冰山模式——\imp{死亡 : 伤残 : 受伤 = 1 : 50 : 1000}。即每发生1例因校园暴力导致的死亡，背后约有50例伤残和1000例受伤事件。校园暴力和人类许多行为不同，它并不伴随社会文明的提高而减少，反而表现为严重化。

\textbf{2. 校园暴力社会生态系统模型：}
校园暴力的发生是社会生态系统多层面因素交互作用的结果：

\begin{enumerate}[leftmargin=*]
    \item \imp{社会因素（宏观层面）}：
    \begin{itemize}
        \item 社会文化：暴力文化的渗透。
        \item 政策和教育体制的不完善。
        \item 媒体和社交网络的不良影响。
        \item 经济与社会不公平。
        \item 社会认知与教育的偏差。
    \end{itemize}
    \item \imp{群体因素（中观层面）}：
    \begin{itemize}
        \item 同伴影响和同辈压力。
        \item 群体排斥和社会隔离。
        \item 集体氛围和文化。
        \item 群体中的竞争和排他性。
        \item 群体行为去个性化效应。
    \end{itemize}
    \item \imp{个体因素（微观层面）}：
    \begin{itemize}
        \item 性格特征：攻击性、冲动性、自尊心和自卑感、社会适应能力差等。
        \item 情绪管理问题和心理创伤。
        \item 暴力认知因素。
        \item 社交经验不足、模仿行为。
        \item 学习和认知发展因素。
    \end{itemize}
\end{enumerate}
\end{keybox}

\subsection{网络暴力}

\begin{keybox}
\textbf{一、定义}
\begin{enumerate}[leftmargin=*]
    \item \imp{网络暴力（cyber violence）}：以互联网给为载体，通过文字、图片、视频等形式对他人实施侮辱、诽谤、威胁、骚扰、侵犯隐私等行为，给受害人造成心理压力和精神损害的暴力形式。
\end{enumerate}

\textbf{二、四大特征}
\begin{enumerate}[leftmargin=*]
    \item \imp{强制性}：施暴者通过反复的信息轰炸、威胁、羞辱等手段，对受害人实施强制性的精神折磨。
    \item \imp{盲目性}：网络施暴者常对受害人的实际情况缺乏了解，仅凭片面信息或群体情绪进行攻击，具有明显的盲目跟风特点。
    \item \imp{匿名性}：网络的匿名性降低了施暴者的责任感和内疚感，易于产生"去抑制效应"，使施暴行为更加肆无忌惮。
    \item \imp{伤害的持久性}：网络信息可被快速、大范围传播且难以彻底删除，对受害者的负面影响持续时间长、范围广。
\end{enumerate}

\textbf{三、主要表现形式}
\begin{enumerate}[leftmargin=*]
    \item \imp{传播负面信息}：在网上散布虚假信息、造谣诽谤、公开羞辱等。
    \item \imp{人肉搜索与公开个人信息（doxxing）}：未经本人同意，在网上公开个人隐私信息（真实姓名、家庭住址、手机号、照片等），使其遭受骚扰和威胁。
    \item \imp{线下延伸}：网络暴力可能延伸到现实中，对受害人的学习和生活造成实际的干扰和伤害。
\end{enumerate}

\textbf{四、危险因素（多层面分析）}
\begin{enumerate}[leftmargin=*]
    \item \imp{微观层面}：个体心理特质（冲动、低共情、低自尊）、网络使用频率和时间、受害经历。
    \item \imp{中观层面}：同伴影响、家庭监管缺失、学校网络素养教育不足。
    \item \imp{宏观层面}：网络监管法律不健全、网络道德规范缺失、社会价值观多元化冲突。
\end{enumerate}

\textbf{五、危害（5项）}
\begin{enumerate}[leftmargin=*]
    \item \imp{加剧青少年的暴力倾向，弱化道德法治观念}：长期接触网络暴力内容可使青少年对暴力行为脱敏，降低对暴力的道德评判和法治敬畏。
    \item \imp{影响青少年认知和人格的健全发展}：网络暴力中的负面信息和扭曲价值观可干扰青少年正常的价值观形成和人格塑造。
    \item \imp{不利于青少年人际沟通，影响身心健康}：网络暴力的受害者和施暴者均可能出现人际交往困难、社交焦虑等问题，对身心健康产生负面影响。
    \item \imp{扰乱虚实社会秩序}：网络暴力破坏了线上线下的正常社会交往秩序，模糊了虚拟与现实的边界，使社会规范在双重空间中都遭到挑战。
    \item \imp{践踏现实社会法律}：严重的网络暴力行为（如人肉搜索、网络诽谤、网络骚扰等）可能构成违法犯罪，侵犯公民合法权益，破坏法治社会基础。
\end{enumerate}

\textbf{六、多层面预防策略}
\begin{enumerate}[leftmargin=*]
    \item \imp{政府层面}：(1)完善网络空间监管和治理的法律法规；(2)明确网络暴力的法律界定和惩治标准；(3)建立网络暴力的举报和处理机制；(4)推进网络实名制和信息溯源技术应用。
    \item \imp{媒体层面}：(1)主流媒体加强正面引导，倡导文明上网；(2)平台企业加强内容审核和管理责任；(3)开发和完善网络暴力的自动识别和过滤技术。
    \item \imp{学校层面}：(1)将网络素养和数字公民教育纳入学校课程；(2)开展反网络暴力专题教育；(3)建立学生网络行为规范和监督机制；(4)对网络暴力事件及时发现、干预和处理。
    \item \imp{家庭层面}：(1)家长提高自身网络素养，以身作则；(2)关注子女网络使用情况（使用时间、内容、行为）；(3)建立良好的亲子沟通，使子女遇到网络暴力时敢于求助；(4)营造温暖和支持性的家庭氛围。
\end{enumerate}
\end{keybox}

\newpage
